\documentclass[11pt]{article}
\usepackage[T1]{fontenc}
\usepackage[utf8]{inputenc}
\usepackage[french]{babel}
\usepackage{newpxtext,newpxmath}
\usepackage[a4paper,top=1.6cm,bottom=1.55cm,left=2.0cm,right=2.0cm]{geometry}
\usepackage{microtype}
\usepackage{xcolor}
\usepackage{titlesec}
\usepackage[most]{tcolorbox}
\usepackage{fancyhdr}
\usepackage{enumitem}
\usepackage{array}
\usepackage{booktabs}
\usepackage{longtable}

\definecolor{efenavy}{HTML}{1F3A5F}
\definecolor{efeteal}{HTML}{0E7A6E}
\definecolor{efeamber}{HTML}{B25A1E}
\definecolor{efegray}{HTML}{54606E}
\definecolor{efesolve}{HTML}{2C6E49}
\definecolor{efered}{HTML}{A32E2E}
\definecolor{boxgray}{HTML}{F3F5F8}
\definecolor{boxamber}{HTML}{FBF3EA}
\definecolor{boxedge}{HTML}{D9E0EA}

\titleformat{\section}{\color{efenavy}\large\bfseries}{}{0pt}{}
\titlespacing{\section}{0pt}{9pt}{3pt}
\titleformat{\subsection}[runin]{\color{efenavy}\bfseries}{}{0pt}{}[\;---\;]
\titlespacing{\subsection}{0pt}{5pt}{0.5em}

\setlength{\parskip}{2pt}
\setlength{\parindent}{0pt}
\setlength{\emergencystretch}{3em}

\pagestyle{fancy}\fancyhf{}
\renewcommand{\headrulewidth}{0.3pt}
\lhead{\small\color{efegray}Scella --- la feuille de route du MVP (C1 $+$ C3)}
\cfoot{\small\color{efegray}\thepage}

\newcommand{\code}[1]{\texttt{\small #1}}
\newcommand{\best}[1]{{\color{efesolve}\textbf{#1}}}
\newcommand{\pire}[1]{{\color{efered}\textbf{#1}}}
\newcommand{\moyen}[1]{{\color{efeamber}\textbf{#1}}}
\newcommand{\mm}[1]{\textbf{#1}}
\newcommand{\sz}[1]{{\small #1}}
\newcommand{\mesure}[1]{{\color{efeteal}\textbf{Mesuré.} }#1\par}
\newcommand{\refute}[1]{{\color{efered}\textbf{Réfuté.} }#1\par}
\newcommand{\fl}{$\rightarrow$}

\newtcolorbox{these}{colback=boxgray,colframe=boxedge,boxrule=0.5pt,arc=2.5pt,
  left=8pt,right=8pt,top=6pt,bottom=6pt}
\newtcolorbox{cadre}{colback=boxamber,colframe=efeamber!55,boxrule=0.6pt,arc=2.5pt,
  left=8pt,right=8pt,top=6pt,bottom=6pt}

\begin{document}

\begin{center}
{\Huge\bfseries\color{efenavy} Le MVP}\\[3pt]
{\large\color{efegray}\itshape ce qui sépare le dépôt d'un produit qu'un inconnu peut lancer
--- C1 $+$ C3, 2 septembre 2026}
\end{center}

\vspace{2pt}

\begin{these}
\textbf{Les dix-sept étapes du programme sont exécutées.} Ce document ne les rejoue pas. Il pose une
question différente, et plus dure : \emph{qu'est-ce qui empêche aujourd'hui une personne extérieure
d'installer cette chose, de la faire tourner, et d'en tirer quelque chose ?}

La réponse tient en quatre points, tous mesurés le 2 septembre 2026, et le premier est le seul qui
compte vraiment : \textbf{la porte d'entrée n'a jamais été ouverte}. Le \code{docker compose up} qui
constitue le premier geste d'un prospect n'a jamais été exécuté --- pas une fois, par personne.
\end{these}

\section{1 \;\textbar\; Ce que « MVP » veut dire ici, et ce qu'il ne veut pas dire}

Un programme achevé et un produit minimal viable ne sont pas la même chose, et confondre les deux est
la façon la plus courante de rester six mois de plus dans un dépôt. Le programme demandait : \emph{les
mécanismes existent-ils, et tiennent-ils sous la mesure ?} La réponse est oui, module par module,
contrôle par contrôle. Le MVP demande autre chose : \emph{le plus petit assemblage qu'un inconnu peut
lancer, qui délivre la promesse du produit, et derrière lequel on peut se tenir sans nuancer.}

La différence n'est pas de degré. Un mécanisme qui passe ses tests unitaires et un service qui démarre
dans un conteneur sur la machine de quelqu'un d'autre sont séparés par une classe entière de
défaillances : le contexte de construction, les couches d'image, les volumes partagés, le réseau, les
variables absentes, la lenteur du premier démarrage. Aucune de ces défaillances n'apparaît dans un
\code{cargo test}. Toutes apparaissent dans les trente premières secondes d'un prospect.

C'est pourquoi ce document ne classe pas les travaux par difficulté ni par élégance, mais par
\textbf{ordre de découverte} : ce qu'un inconnu rencontre en premier passe en premier, même si c'est
techniquement le moins intéressant.

\section{2 \;\textbar\; Le fait qui décide de la forme du MVP}

Avant toute planification, une mesure change la question. On parle de « C1 $+$ C3 » comme de deux
produits à assembler. Ils n'en sont pas.

\begin{center}\small
\renewcommand{\arraystretch}{1.2}
\begin{tabular}{@{}l l r@{}}
\toprule
\textbf{Binaire} & \textbf{Features} & \textbf{Crates du dépôt} \\
\midrule
\code{efe-swarm-gateway} & défaut & 13 \\
\code{efe-swarm-gateway} & $+$ \code{reg-evidence} & \best{14} \\
\code{audit-console} & défaut & 7 \\
\bottomrule
\end{tabular}
\end{center}

\mesure{\code{cargo tree --edges normal}, crates du dépôt uniquement, le 2 septembre 2026.}

\textbf{Allumer C1 sur un C3 déployé coûte exactement un crate.} Ce n'est pas une commodité
d'ingénierie, c'est la thèse commerciale rendue littérale : le registre réglementaire n'est pas un
second produit qu'il faudrait vendre et installer, c'est un \emph{interrupteur} sur un déployable que
le client fait déjà tourner. La feuille de route de \mm{M5} avait posé ce principe --- « joindre C1 et
C3 \emph{sans les souder} », feature éteinte par défaut, garde qui lit le manifeste --- et la mesure
confirme que le principe a survécu à son implémentation.

La conséquence pour le MVP est directe et simplificatrice : \textbf{il n'y a pas deux installations à
concevoir}. Il y en a une, et une case à cocher. Tout plan qui traiterait C1 comme un déploiement
séparé ajouterait un travail que le code ne demande pas.

\section{3 \;\textbar\; Ce qui sépare aujourd'hui du MVP --- quatre faits mesurés}

\subsection{(a) La porte d'entrée n'a jamais été ouverte}
Le dépôt contient \code{compose.yaml} (deux services : \code{passerelle} et \code{preuve}) et
\code{docker/Dockerfile.coin}, une image \emph{distroless} portant trois binaires --- la passerelle,
le pilote de démonstration, et \code{verify-mini} dont l'arbre de dépendances est vide. La
configuration est valide : \code{docker compose config} rend les deux services.

Mais \code{m4-check.sh} le dit lui-même, en toutes lettres, dans son propre en-tête : le critère de
\mm{M4} \emph{se sépare en deux moitiés} --- la charge utile, qui s'exécute sans conteneur, et
l'emballage, qui exige un démon Docker --- et « \emph{ce script prouve la première et ne prétend rien
de la seconde} ». La seconde moitié n'a jamais tourné, faute de démon.

\mesure{Le démon est désormais présent : Docker \code{29.4.0}, Compose \code{v5.1.1}. La contrainte
qui rendait cette vérification impossible est levée.}

\subsection{(b) La porte de qualité du dépôt est rouge, et elle ne s'ouvre jamais}
\refute{\code{cargo clippy -{}-workspace -{}-all-targets -{}- -D warnings} --- la commande
\emph{exacte} de \code{ci.yml:40}, le job \textbf{bloquant} --- rend \pire{exit 101} :
\textbf{109 erreurs sur 6 crates} (\code{efe-events} 55, \code{efe-core} 26, \code{efe-prompts} 19,
\code{efe-nano-app-coordinator} 3, \code{verify-meta-mini} 1, \code{efe-kg} 1).}

Deux choses aggravent ce chiffre. D'abord, personne ne pouvait le voir : les déclencheurs du workflow
sont \code{push} sur \code{main} et \code{release/**}, et \code{pull\_request} vers \code{main} ---
or tout le travail vit sur \code{couche2-m0-frontier}. La « CI verte » n'était pas verte, elle était
\textbf{absente}. Ensuite, le relevé que le dépôt tient par écrit --- environ 53 sites sur 7 crates,
inscrit au-dessus du job \code{clippy-frontiere}, qui est \emph{non bloquant} --- décrit une dette
\textbf{différente} : un seul de ces sept crates figure parmi les six ci-dessus. Environ 106 erreurs
de la porte bloquante ne sont écrites nulle part.

Il faut être honnête sur ce que cela pèse : un prospect ne regarde pas les lints. Ce point ne bloque
donc pas l'usage, il bloque le \emph{discours}. On ne peut pas dire « le dépôt tient ses propres
portes » d'une porte qui n'a jamais été franchie.

\subsection{(c) La suite complète n'a jamais été vérifiée verte de bout en bout}
\mesure{Le 2 septembre, \code{cargo test -{}-workspace} rend \textbf{11\,235 verts et un seul échec}
--- une course sur \code{std::env} dans \code{efe-nano-inference-daemon}, désormais supprimée par
construction (15/20 \fl{} \best{30/30} en rejeux parallèles).}

Ce qui manque est la \emph{confirmation} : le rejeu complet après correction n'a pas pu se faire, la
machine paginant lourdement (13,4~Gio de swap sur 14,3). Sous pagination, un tour de workspace se
compte en heures --- et une mesure qu'on ne peut pas croire ne vaut pas mieux qu'une absence de
mesure. Le critère de \mm{R3} en demande d'ailleurs \emph{trois} d'affilée.

\subsection{(d) Le registre n'est pas dans la démonstration}
La feature \code{reg-evidence} est éteinte par défaut --- délibérément, et \mm{M5} explique pourquoi.
Mais \code{compose.yaml} construit l'image avec les features par défaut : \textbf{le
\code{docker compose up} d'aujourd'hui montre C3 et ne montre pas C1}. Pour un MVP annoncé
« C1 $+$ C3 », c'est la moitié du produit qui reste invisible au premier contact.

\section{4 \;\textbar\; Les étapes, dans l'ordre, avec leur critère}

\begin{center}\small
\renewcommand{\arraystretch}{1.3}
\begin{tabular}{@{}>{\raggedright\arraybackslash}p{0.7cm} >{\raggedright\arraybackslash}p{4.1cm} >{\raggedright\arraybackslash}p{8.3cm} >{\raggedright\arraybackslash}p{2.1cm}@{}}
\toprule
& \textbf{Étape} & \textbf{Le critère --- ce qu'on doit pouvoir exécuter} & \textbf{Coût} \\
\midrule
\mm{É1} & \textbf{Ouvrir la porte} & \code{docker compose up} sur cette machine : les deux services
démarrent, \code{preuve} sort 0, un reçu est exporté et re-prouvé \textbf{sans clé}, la troncature est
démasquée & \sz{1--2 j} \\
\mm{É2} & \textbf{Allumer C1 dans la \emph{même} commande} & le compose construit avec
\code{reg-evidence} ; la démonstration dérive un dossier réglementaire \textbf{du même run}, et
\code{efe\_swarm\_reg\_pont} vaut 1 & \sz{1--2 j} \\
\mm{É3} & \textbf{La suite, confirmée} & \code{r3-check.sh} \fl{} 0 avec \code{R3\_TOURS=3}, sur une
machine qui ne pagine pas & \sz{$\frac{1}{2}$ j} \\
\mm{É4} & \textbf{La porte de qualité, verte \emph{et} déclenchée} & \code{clippy -{}-workspace
-{}-all-targets -D warnings} \fl{} 0, \textbf{et} le workflow se déclenche sur la branche de travail
& \sz{3--5 j} \\
\mm{É5} & \textbf{La page qu'on tend} & une page : la commande, ce qu'on va voir, ce qu'on peut
vérifier soi-même, et ce que la démonstration \textbf{ne prouve pas} & \sz{1 j} \\
\bottomrule
\end{tabular}
\end{center}

\subsection{Pourquoi \mm{É1} passe avant \mm{É4}}
L'ordre inverse serait plus flatteur : réparer sa maison avant d'ouvrir. Il serait faux ici, pour une
raison précise. Les 109 erreurs de la porte bloquante sont des \emph{lints} --- documentation
manquante, signatures trop longues, motif remplaçable par \code{clamp} --- et non des défauts de
correction : \code{cargo check -{}-workspace -{}-all-targets} passe. Elles ne changent donc rien à ce
qu'un prospect voit, ni à ce que le produit fait. Tandis que \mm{É1} peut révéler une classe entière
de défaillances qu'aucun test du dépôt n'a jamais pu atteindre, et dont on ne saura rien tant qu'on
n'aura pas essayé.

Il y a un argument plus fort encore : \mm{É1} est le seul point du plan dont on ignore réellement le
coût. Les 109 erreurs se comptent et se corrigent ; le premier build de \code{Dockerfile.coin}
compile la fermeture de la passerelle depuis zéro dans un conteneur, et son en-tête prévient qu'il est
\emph{long}. \textbf{On planifie mal autour d'une inconnue qu'on garde pour la fin.}

\subsection{Ce qu'il faut attendre de \mm{É1}}
Non pas un succès, mais une \emph{information}. Un chemin de construction jamais exercé échoue
rarement du premier coup, et les défaillances probables sont connues d'avance : un fichier absent du
contexte de build, un binaire non déclaré dans l'étape finale, un volume monté à un chemin que
l'image distroless ne possède pas, un profil dont une variable manque en conteneur. Aucune n'est
grave ; toutes coûtent une demi-journée chacune si on les découvre une par une devant un prospect.

\section{5 \;\textbar\; Ce que le MVP n'est pas, et pourquoi le dire ici}

Un plan qui ne dit pas ce qu'il exclut se fait élargir par tout le monde. Les exclusions ci-dessous
sont \emph{déjà} des décisions du dépôt ; les répéter évite qu'on les rouvre par inadvertance.

\begin{itemize}[leftmargin=1.1em,itemsep=1.5pt,topsep=2pt]
\item \textbf{Le collecteur ouvert (\mm{MX}) reste hors du MVP.} La consigne écrite est catégorique :
n'écrire aucune ligne avant que trois personnes aient payé pour l'obtenir. Le compteur est à
\pire{zéro}.
\item \textbf{Le processus ne termine pas TLS.} Il exige la \emph{déclaration} du terminateur placé
devant lui, refuse de servir hors boucle locale sans elle, et depuis le 2 septembre il peut
\emph{corroborer} cette déclaration par l'adresse source. Tirer une pile TLS dans un crate feuille
serait moins sûr qu'un terminateur mûr.
\item \textbf{Le jeton porteur reste un secret partagé.} Une vraie non-répudiation demanderait que
l'officier \emph{signe} le tuple de l'acte avec une clé que le serveur ne détient pas. C'est nommé,
et non produit.
\item \textbf{Le barreau de canal n'est pas scellé dans le nœud \code{Human}.} Le faire changerait le
format, ce qui est un \emph{événement de version}, pas une étape de MVP. Le service \textbf{coupe} au
lieu d'enregistrer : un acte qui ne franchit pas le plancher n'est jamais scellé.
\end{itemize}

\section{6 \;\textbar\; Le jalon qui ne se code pas, et qui décide seul}

\begin{cadre}
Une fois \mm{É1} à \mm{É5} faites, il ne restera aucune raison technique d'attendre. Le mode
observation existe depuis \mm{M2}, le paquetage depuis \mm{M4}, la sauvegarde et la rotation depuis
\mm{P2}, l'ancrage tiers depuis \mm{M9}. Ce qui manquera alors n'est pas dans ce document :
\textbf{aucune cible extérieure n'est déclarée depuis le 11 août 2026}.

Le seul jalon qui décide reste donc : \emph{un prospect installe-t-il C3 en mode observation ?} S'il
est atteint et que personne n'installe, le problème n'est pas le produit --- c'est la distribution, et
aucune des cinq étapes ci-dessus n'y changera quoi que ce soit. Le bon moment pour le découvrir est
\textbf{après \mm{É2}}, pas après \mm{É5} : dès que la démonstration montre les deux moitiés du
produit, elle est montrable.
\end{cadre}

\vspace{4pt}
{\color{efegray}\hrule height 0.4pt}
\vspace{2.5pt}
{\footnotesize\color{efegray}Toutes les mesures de ce document ont été produites le 2 septembre 2026
sur le dépôt à jour (\code{64f519a}). Les contrôles se rejouent par
\code{bash scripts/<nom>-check.sh} ; les chiffres de fermeture par
\code{cargo tree -p <binaire> -{}-edges normal}.}

\end{document}
